\documentclass[journal]{IEEEtran}

% ---------------------------------------------------------------------------
%  BLG 475E Software Quality and Testing - Term Project
%  Phase 2 Report (IEEE journal style)
%  Authors: Mustafa Eren KOC (150190805), Onat Baris ERCAN (150210075)
% ---------------------------------------------------------------------------

\IEEEoverridecommandlockouts

\usepackage{cite}
\usepackage{amsmath,amssymb,amsfonts}
\usepackage{algorithm}
\usepackage{algorithmic}
\usepackage{graphicx}
\usepackage{textcomp}
\usepackage{xcolor}
\usepackage{booktabs}
\usepackage{array}
\usepackage{url}
\usepackage{hyperref}
\usepackage{listings}

\hypersetup{
  colorlinks=true,
  linkcolor=black,
  citecolor=black,
  urlcolor=blue
}

\lstdefinestyle{javamini}{
  basicstyle=\ttfamily\scriptsize,
  breaklines=true,
  columns=fullflexible,
  keepspaces=true,
  language=Java,
  xleftmargin=0pt,
  aboveskip=0pt,
  belowskip=0pt,
  frame=none,
  showstringspaces=false
}

\def\BibTeX{{\rm B\kern-.05em{\sc i\kern-.025em b}\kern-.08em
    T\kern-.1667em\lower.7ex\hbox{E}\kern-.125emX}}

\begin{document}

\title{LLM-Based Integration Testing for a Composite Java \texttt{BookScan}
Class:\\ Prompt Combination, Prompt Refinement, and Coverage-Guided Validation}

\author{\IEEEauthorblockN{Mustafa Eren Ko\c{c}}
\IEEEauthorblockA{\textit{Department of Computer Engineering} \\
\textit{Istanbul Technical University}\\
Istanbul, T\"urkiye \\
150190805 \\
\texttt{koc19@itu.edu.tr}}
\and
\IEEEauthorblockN{Onat Bar{\i}\c{s} Ercan}
\IEEEauthorblockA{\textit{Department of Computer Engineering} \\
\textit{Istanbul Technical University}\\
Istanbul, T\"urkiye \\
150210075 \\
\texttt{ercano21@itu.edu.tr}}
}

\maketitle

% ---------------------------------------------------------------------------
\begin{abstract}
Phase~2 of the BLG~475E term project extends Phase~1 from isolated
HumanEval-X Java functions to a single integrated class, \texttt{BookScan},
that must combine the semantics of HumanEval-X tasks \texttt{Java/18}
(substring counting), \texttt{Java/23} (string length) and \texttt{Java/27}
(upper-/lower-case conversion) into one coherent API. The class determines
how many words of a given length appear in a text, on which lines they
appear, and how many times a requested whole word occurs across lines. We
evaluate two small open-weight coder models, \texttt{Qwen2.5-Coder-1.5B} and
\texttt{DeepSeek-Coder-1.3B}, under two prompt strategies (an \emph{original
combined prompt} and an \emph{edited combined prompt}) and compare both with
a manually selected final implementation. A custom prompt-comparison harness
drives \texttt{javac}, JUnit~6 and the JaCoCo CLI to measure compile success,
JUnit pass rate and branch coverage for all five configurations, and the
final implementation is additionally validated through a Maven build that
runs Surefire and produces a JaCoCo report. The original combined prompt
leaves several integration contracts implicit and both models therefore fail
multiple test cases (Qwen $6/18$, DeepSeek $8/18$); the edited combined
prompt makes token boundaries, line-number semantics, whole-word matching
and invalid-input handling explicit and lifts both models to $17/18$ under
the strengthened suite. The selected final implementation closes the
remaining private-helper normalization gap, passes all $18$ JUnit tests,
reaches $72/72 = 100.00\,\%$ branch coverage and $110/110 = 100.00\,\%$
line coverage with JaCoCo, and records a PITest mutation score of
$79/93 = 84.95\,\%$. Our results indicate that at the integration level
the prompt specification matters at least as much as the choice of
underlying model.
\end{abstract}

\begin{IEEEkeywords}
Large language models, integration testing, JUnit, JaCoCo, branch coverage,
mutation testing, PITest, prompt engineering, HumanEval-X, software quality.
\end{IEEEkeywords}

% ---------------------------------------------------------------------------
\section{Introduction}\label{sec:intro}

Phase~1 of the project studied how two open-weight code-LLMs behave on
\emph{isolated} HumanEval-X Java tasks~\cite{humaneval,phase1report}. That
setting is attractive because each task can be specified by a single
function signature plus a few examples, but it does not capture the
difficulty of integrating several helpers into one class that shares data
structures and normalization rules.

Phase~2 moves to exactly that setting. We define a custom class,
\texttt{BookScan}, that must act as if the three HumanEval-X tasks
\texttt{Java/18} (substring counting), \texttt{Java/23} (string length) and
\texttt{Java/27} (upper-/lower-case conversion) were always intended to be
used together. The class is required to:

\begin{itemize}
  \item count how many words of a requested length appear in a text and on
    which lines,
  \item count how many times a requested whole word appears across lines,
  \item preserve the Phase~1 helper semantics of
    \texttt{howManyTimes}, \texttt{strlen} and \texttt{flipCase} so that
    they remain callable and individually testable.
\end{itemize}

We then ask a specific question: \emph{once the prompt has to specify a
whole class rather than a single function, which matters more -- the
underlying model or the structure of the prompt itself?}  The experimental
apparatus that we build to answer that question is a prompt-comparison
harness that runs the same JUnit suite against four candidate
\texttt{BookScan} implementations (two prompt strategies $\times$ two
models) and against a manually selected final implementation. Coverage is
measured uniformly through a scripted JaCoCo CLI pipeline and, for the
final implementation, verified a second time through Maven.

% ---------------------------------------------------------------------------
\section{Team Composition and Carry-over From Phase~1}\label{sec:team}

As reported in Phase~1, our group ended with two members after one teammate
withdrew early in the semester. Phase~2 preserves the same two-person
working structure and the same two models,
\texttt{Qwen2.5-Coder-1.5B-Instruct-GGUF} and
\texttt{DeepSeek-Coder-1.3B-Instruct-GGUF}, so that Phase~1 and Phase~2
conclusions remain comparable. Detailed role allocation is summarized in the
Acknowledgement.

% ---------------------------------------------------------------------------
\section{Experimental Setup}

\subsection{Models}

The two models from Phase~1 are reused unchanged:

\begin{itemize}
  \item \texttt{Qwen/Qwen2.5-Coder-1.5B-Instruct} (GGUF build),
  \item \texttt{deepseek-ai/DeepSeek-Coder-1.3B-Instruct} (GGUF build).
\end{itemize}

Both are executed through the same local \texttt{llama.cpp} runtime with
identical sampling settings. Keeping the model set fixed isolates the
effect of the prompt strategy.

\subsection{Prompt Strategies}\label{sec:prompts}

Two prompt families are compared. The original combined prompt simply
merges the three HumanEval-X task descriptions and asks the model to build
a class around them. The edited combined prompt adds a short list of
integration contracts that the original prompt leaves implicit
(Section~\ref{sec:prompt-diff}).

\subsubsection{Original Combined Prompt}

\begin{quote}\itshape\small
Create a Java class named \texttt{BookScan}. The class should determine how
many times words of a given length appear in a text and on which lines they
appear. The class must also contain methods related to HumanEval-X tasks
\texttt{\#18 substring}, \texttt{\#23 string length}, and \texttt{\#27
upper-lower case}. Return a single compilable Java file. Keep the
implementation simple and self-contained.
\end{quote}

\subsubsection{Edited Combined Prompt}

The edited prompt adds the integration contracts shown in
Table~\ref{tab:edits}.

\begin{table}[h]
\caption{Contract clarifications added in the edited prompt}
\label{tab:edits}
\centering
\begin{tabular}{@{}p{2.6cm}p{4.9cm}@{}}
\toprule
\textbf{Topic} & \textbf{Added requirement} \\
\midrule
Line numbering & one-based, each matching line listed at most once \\
Tokenization   & apostrophes and hyphens stay inside a single token \\
Whole-word search & must not match inside a longer word \\
Case policy    & search is case-insensitive, uppercase-only queries canonicalized \\
Query hygiene  & trim the query, reject blank or multi-token queries \\
Invalid inputs & null or blank text returns empty-result records \\
Result shape   & records \texttt{ScanResult}/\texttt{WordSearchResult}, immutable collections \\
Helper reuse   & \texttt{howManyTimes}/\texttt{strlen}/\texttt{flipCase} remain callable with Phase~1 semantics \\
\bottomrule
\end{tabular}
\end{table}

\subsection{Toolchain}

The Phase~2 pipeline is built around the same tool stack as Phase~1 and
extends it with a Maven build:

\begin{itemize}
  \item JDK 21 (Temurin),
  \item Maven 3.9+ with Surefire 3.5 and JaCoCo 0.8.12,
  \item JUnit~6 (\texttt{junit-platform-console-standalone-6.0.3}),
  \item PITest~1.20.5 plus \texttt{pitest-junit5-plugin~1.2.3} exposed
    through an opt-in Maven profile \texttt{mutation}
    (\texttt{mvn -P mutation test}) so that \texttt{BookScan}'s mutation
    score can be measured automatically without slowing down the default
    \texttt{mvn clean verify} cycle,
  \item PowerShell~7 + Python~3.11 driver scripts for the prompt-comparison
    harness,
  \item GitHub Actions workflow
    \texttt{.github/workflows/phase2-ci.yml} that runs \texttt{mvn clean
    verify}, executes the prompt-comparison harness and -- as an opt-in
    best-effort step -- invokes the \texttt{mutation} profile to upload
    PITest's HTML/XML/CSV reports as build artefacts.
\end{itemize}

\subsection{Prompt-Comparison Harness}

The Python runner \texttt{scripts/python/run\_phase2\_prompt\_comparison.py}
is invoked through the PowerShell wrapper
\texttt{scripts/Run-Phase2PromptComparison.ps1}. For every candidate
\texttt{BookScan.java} (four prompt variants plus the manually selected
final implementation) it:

\begin{enumerate}
  \item copies the candidate and the shared JUnit suite into an isolated
    working directory;
  \item compiles with \texttt{javac} against the JUnit~6 console jar;
  \item runs JUnit under the JaCoCo agent;
  \item renders a JaCoCo CSV from the resulting \texttt{jacoco.exec};
  \item extracts branch coverage for the \texttt{BookScan} class.
\end{enumerate}

All raw artefacts (\texttt{compile.stdout}, \texttt{compile.stderr},
\texttt{junit.stdout}, \texttt{jacoco.exec}, \texttt{jacoco.csv} and the
copied source files) are stored under
\texttt{Phase2/results/artifacts/prompt-comparison/\{strategy\}/\{model\}}
so that every figure in this report can be reproduced from recorded data
rather than from in-memory state.

% ---------------------------------------------------------------------------
\section{\texttt{BookScan} Design}\label{sec:design}

The final \texttt{BookScan} class exposes two public record results and
five public methods (Fig.~\ref{fig:api}). The design follows three
principles that the Phase~1 experience highlighted as important:

\begin{itemize}
  \item \textbf{Make every helper individually callable} so Phase~1
    regression tests can keep running.
  \item \textbf{Use immutable records} at the public boundary so
    callers cannot mutate scan results.
  \item \textbf{Normalize inputs once} inside the public entry points so
    defensive guards do not spread across the private helpers.
\end{itemize}

\begin{figure}[h]
  \centering
  \framebox[\columnwidth]{\parbox[c][4.5cm][c]{0.94\columnwidth}{
    \centering
    \textbf{Public API of} \texttt{BookScan} \\[2pt]
    \texttt{record ScanResult(targetLength, totalOccurrences,}\\
    \texttt{~~matchingLines, occurrencesByLine, matchedWordsByLine)} \\
    \texttt{record WordSearchResult(word, totalOccurrences,}\\
    \texttt{~~matchingLines, occurrencesByLine)} \\[2pt]
    \texttt{ScanResult scanByWordLength(text, targetLength)} \\
    \texttt{WordSearchResult scanWord(text, word)} \\
    \texttt{int howManyTimes(string, substring)} \\
    \texttt{int strlen(string)} \\
    \texttt{String flipCase(string)}
  }}
  \caption{Public surface of the final \texttt{BookScan} class.}
  \label{fig:api}
\end{figure}

Tokenization uses the compiled pattern

\begin{center}
\texttt{[A-Za-z0-9]+(?:['-][A-Za-z0-9]+)*}
\end{center}

which keeps apostrophes and hyphens inside a token (\texttt{can't},
\texttt{co-op}) while still rejecting them as leading or trailing
characters. Whole-word matching is implemented by canonicalizing each line
into a space-separated list of lowercase tokens and then invoking
\texttt{howManyTimes} on a padded needle so that partial matches inside
longer words (e.g.\ \texttt{echo} inside \texttt{echoic}) cannot fire.

% ---------------------------------------------------------------------------
\section{Integration Test Strategy}\label{sec:tests}

The Phase~2 JUnit suite contains $12$ integration tests and $6$ focused
regression / hardening tests, for a total of $18$ tests shared by all five
configurations. The integration tests are summarized in
Table~\ref{tab:int-tests} and the regression tests in
Table~\ref{tab:reg-tests}.

\begin{table}[h]
\caption{Integration tests (\texttt{BookScanIntegrationTest})}
\label{tab:int-tests}
\centering
\begin{tabular}{@{}p{3.2cm}p{4.3cm}@{}}
\toprule
\textbf{Test method} & \textbf{Integration contract} \\
\midrule
\texttt{scanByWordLengthAggregatesCountsAndLinesAcrossText} & multi-line aggregation and per-line matched words \\
\texttt{scanByWordLengthTreatsApostrophesAndHyphensAsSingleWords} & token-boundary preservation \\
\texttt{scanByWordLengthReturnsEmptyResultForInvalidInputs} & null / blank / non-positive length \\
\texttt{scanByWordLengthKeepsMatchingLinesUniqueEvenWithManyHits} & line-list uniqueness under repeated hits \\
\texttt{scanWordCountsWholeWordMatchesIgnoringCaseAcrossLines} & case-insensitive whole-word matching \\
\texttt{scanWordFindsTrimmedQueriesNextToPunctuation} & trimmed queries, punctuation handling \\
\texttt{scanWordRejectsBlankOrMultiTokenQueries} & query validation (blank, multi-token) \\
\texttt{scanWordReturnsEmptyResultForNullInputsAndBlankTexts} & null / blank inputs produce stable empty results \\
\texttt{scanWordDoesNotMatchInsideLongerWords} & partial-match rejection \\
\texttt{scanWordCanonicalizesUppercaseQueriesBeforeWholeWordMatching} & uppercase queries normalized via \texttt{flipCase} \\
\texttt{scanWordHandlesAlphanumericWholeWordQueries} & alphanumeric whole-word tokens \\
\texttt{scanWordNormalizesMixedCaseTokensWithoutFlipArtifacts} & mixed-case normalization remains stable \\
\bottomrule
\end{tabular}
\end{table}

\begin{table}[h]
\caption{Regression tests (\texttt{BookScanRegressionTest})}
\label{tab:reg-tests}
\centering
\begin{tabular}{@{}p{3.2cm}p{4.3cm}@{}}
\toprule
\textbf{Test method} & \textbf{Phase~1 helper under test} \\
\midrule
\texttt{howManyTimesCountsOverlappingSubstrings} & \texttt{Java/18} overlap semantics \\
\texttt{howManyTimesReturnsZeroForEmptyNullOrOversizedNeedles} & \texttt{Java/18} guard on blank / oversized needles \\
\texttt{flipCasePreservesSymbolsWhileTogglingLetters} & \texttt{Java/27} case-flip on symbols and digits \\
\texttt{strlenThrowsOnNullAndReturnsLengthOtherwise} & \texttt{Java/23} null contract and length \\
\texttt{tokenizeWordsReturnsFreshMutableEmptyListForNullInput} & helper hardening for null tokenization \\
\texttt{privateCanonicalizeWordRejectsBlankTokens} & helper hardening for blank-token collapse \\
\bottomrule
\end{tabular}
\end{table}

The integration tests were authored from a cross-product of three sources:
the Phase~1 test-smell checklist, the Phase~2 black-box assessment
(\texttt{docs/analysis/black\_box\_assessment.md}), and the concrete
failure signatures we observed while running the original-prompt variants.
Every test asserts at least one structural property (counts) and at least
one shape property (matching-line list or per-line map) so that a pass
genuinely implies integration correctness, not just one scalar equality.

% ---------------------------------------------------------------------------
\section{Results}\label{sec:results}

\subsection{Prompt-Comparison Results}\label{sec:prompt-diff}

Table~\ref{tab:compare} aggregates the raw CSV in
\texttt{Phase2/results/prompt\_comparison\_results.csv}. The original
prompt lets both models compile but produces significant JUnit failures;
the edited prompt removes all public-integration failures but still leaves
one helper-hardening regression in each model; the selected final
implementation closes that gap and reaches full coverage.

\begin{table}[h]
\caption{Prompt-comparison results. Coverage is reported for the
\texttt{BookScan} class only.}
\label{tab:compare}
\centering
\begin{tabular}{@{}llrr@{}}
\toprule
\textbf{Strategy} & \textbf{Model} & \textbf{Tests passed} & \textbf{Branch cov.} \\
\midrule
original-combined & Qwen     & 6 / 18  & n/a \\
original-combined & DeepSeek & 8 / 18  & n/a \\
edited-combined   & Qwen     & 17 / 18 & n/a \\
edited-combined   & DeepSeek & 17 / 18 & n/a \\
selected-final    & manual   & 18 / 18 & 72 / 72 = 100.00\,\% \\
\bottomrule
\end{tabular}
\end{table}

\subsection{Original Prompt Failure Modes}

The recorded JUnit output classifies the original-prompt failures into
four Qwen defect classes and four DeepSeek defect classes
(Table~\ref{tab:failure-modes}). All of them are integration-level
contracts that are missing from the original prompt and that would have
been invisible at the single-function Phase~1 level.

\begin{table}[h]
\caption{Original combined prompt: defect classes observed through JUnit}
\label{tab:failure-modes}
\centering
\begin{tabular}{@{}p{1.6cm}p{5.9cm}@{}}
\toprule
\textbf{Model} & \textbf{Defect classes hit by the JUnit suite} \\
\midrule
Qwen     & zero-based / non-unique line numbers; raw space-splitting destroys \texttt{can't}/\texttt{co-op}; case-sensitive substring search; no canonicalization for uppercase or alphanumeric queries; missing helper extraction for the final hardening checks \\
DeepSeek & punctuation stripping changes token identity; substring matching inside longer words; over-counting in case-insensitive search; weak handling of punctuation around trimmed queries; alphanumeric over-counting; missing helper extraction for the final hardening checks \\
\bottomrule
\end{tabular}
\end{table}

\subsection{Edited Prompt Improvement}

Both edited-prompt variants clear every public integration scenario, but
under the strengthened suite they stop at $17/18$ because they still fail
\texttt{privateCanonicalizeWordRejectsBlankTokens}. The edited prompt
therefore closes \emph{more} defects than switching between the two
models, while the final manual selection is still needed to harden one
remaining helper contract. This is consistent with the recent literature
on LLM unit-test generation, where explicit specifications have been
repeatedly shown to outperform model upsizing~\cite{metatestgen,yang2024,panta2024,lit_review}.

\subsection{Final Implementation Coverage}\label{sec:final-cov}

The manually selected final implementation is verified twice: by the
prompt-comparison harness and by a full Maven build. The Maven JaCoCo CSV
at \texttt{Phase2/target/site/jacoco/jacoco.csv} reports

\begin{itemize}
  \item branch coverage $72 / 72 = 100.00\,\%$,
  \item line coverage $110 / 110 = 100.00\,\%$,
  \item cyclomatic complexity covered $51 / 51 = 100.00\,\%$.
\end{itemize}

No JaCoCo branch or line misses remain in the final Maven report.

\subsection{Strengthened Regression Closure}\label{sec:residual}

\begin{itemize}
  \item \texttt{tokenizeWordsReturnsFreshMutableEmptyListForNullInput}
    hardens the null-input behavior of the tokenization helper.
  \item \texttt{privateCanonicalizeWordRejectsBlankTokens} pins down the
    blank-token normalization rule that both edited-prompt variants still
    miss.
\end{itemize}

These helper-hardening regressions close the last structural coverage gaps
and make the selected final implementation the only variant that satisfies
the full $18$-test suite.

\subsection{Carry-over From Phase~1 Mutation Guardrails and PITest}\label{sec:mutants}

The Phase~1 task corpus already has $60$ hand-crafted
\texttt{improvedMutation\ldots} JUnit methods, each pinned to a specific
operator family (ROR / AOR / LCR / UOI / SBR / RV / NPE / BND /
EXC)~\cite{phase1report}. Three of those task suites -- \texttt{Java/18},
\texttt{Java/23} and \texttt{Java/27} -- exercise exactly the helpers that
\texttt{BookScan} now reuses, and their semantics are re-verified by the
six Phase~2 regression / hardening tests in Table~\ref{tab:reg-tests}. In other
words, the Phase~1 mutation-based guardrails continue to protect Phase~2
at the helper level, even though Phase~2 now adds two extra hardening
regressions on top of those inherited contracts.

To complement those hand-written guardrails we also wire PITest
(\texttt{org.pitest:pitest-maven}~1.20.5 with
\texttt{pitest-junit5-plugin}~1.2.3) into a Maven profile called
\texttt{mutation}. Running \texttt{mvn -P mutation test} instruments
\texttt{BookScan.class}, re-runs the $18$ integration/regression tests
against each generated mutant and emits an HTML/XML/CSV report under
\texttt{Phase2/target/pit-reports/}. The profile is opt-in so that the
default \texttt{mvn clean verify} remains fast, but the GitHub Actions
workflow invokes it as a best-effort step and uploads the PITest reports
alongside the other Phase~2 artefacts. The combination of PITest
(automated operator coverage on \texttt{BookScan}) and the Phase~1
hand-crafted guardrails (targeted operator coverage on the reused
helpers) gives the Phase~2 submission two independent mutation lenses.
The latest measured PITest result is $79/93 = 84.95\,\%$ killed mutants
with $110/110$ mutated lines covered. The remaining survivors are
concentrated in equivalent or defensive helper mutations rather than in
untested public behaviour.

% ---------------------------------------------------------------------------
\section{Related Work}

Our methodology is most closely aligned with the
validation-plus-repair style of ChatUniTest~\cite{chatunitest} and the
industrial filtering reported by Meta's TestGen-LLM~\cite{metatestgen}.
Yang et al.~\cite{yang2024} and Panta et al.~\cite{panta2024} argue that
prompt structure and static/dynamic feedback loops are more effective than
raw model size, which matches our observation that prompt clarity
outweighs model choice at the integration level. The 2026 systematic
literature review by Tasarsu, Tokmak and Catal~\cite{lit_review} supports
treating branch coverage and behavioural tests as \emph{complementary}
metrics -- exactly the stance we take in
Section~\ref{sec:final-cov}--\ref{sec:mutants}.

% ---------------------------------------------------------------------------
\section{Threats to Validity}

Internal validity is limited by the local-model non-determinism inherent
to \texttt{llama.cpp}. We mitigate this by keeping prompts byte-identical
across runs, pinning the sampling parameters, and storing each
configuration's raw output under
\texttt{Phase2/results/artifacts/prompt-comparison/}. External validity
is limited because Phase~2 evaluates a \emph{single} composite class
rather than an entire project; nevertheless, \texttt{BookScan} is
deliberately constructed to stress three distinct Phase~1 helpers
simultaneously, which is stronger than the single-function HumanEval-X
setting. Construct validity depends on our integration-test suite being a
faithful proxy for correct behaviour; we partially offset this risk
through the black-box assessment
(\texttt{docs/analysis/black\_box\_assessment.md}), which ties every
public equivalence class to at least one executed JUnit method.

% ---------------------------------------------------------------------------
\section{Conclusion and Phase~2 $\to$ Submission Status}

Phase~2 completes the transition from isolated HumanEval-X tasks to a more
realistic integration setting. Both models are capable of producing a
valid \texttt{BookScan} implementation, but only after the prompt
specifies the interaction contracts that the original combined prompt
leaves implicit. The selected final implementation passes the full
$18/18$ suite, reaches $100.00\,\%$ branch coverage and
$100.00\,\%$ line coverage under the Maven build, records an
$84.95\,\%$ PITest mutation score, and preserves all Phase~1 helper-level
mutation guardrails. The repository is therefore ready for submission
and, crucially, the Phase~2 results strengthen rather than contradict the
Phase~1 headline: \emph{once the guardrails are in place, prompt
engineering does more work than model selection}.

% ---------------------------------------------------------------------------
\section*{Acknowledgement}

The source repository, including the prompt-comparison harness, the
Maven build, the PITest profile, the GitHub Actions workflow and the
raw JaCoCo artefacts, is available at
\url{https://github.com/mustafaerenkoc44/BLG-475E-Term-Project}. The
Phase~1 split into engineering and analysis tracks carried over to
Phase~2; both authors contributed equally and cross-reviewed each
other's deliverables.

\textbf{Mustafa Eren KO\c{C}} (150190805) led the
\emph{engineering and automation} track: the Maven and PITest
configuration, the prompt-comparison harness
(\texttt{run\_phase2\_prompt\_comparison.py}), the \texttt{BookScan}
reference implementation and helper hardening, the JaCoCo automation,
the GitHub Actions workflow, and the integration-test driver.

\textbf{Onat Bar{\i}\c{s} ERCAN} (150210075) led the \emph{quality and
analysis} track: the original-versus-edited prompt-strategy design,
the \texttt{BookScan} black-box assessment with per-method equivalence
classes, the failure-mode classification for the original-prompt
variants, the integration-test specification, and the IEEE journal
manuscript.

% ---------------------------------------------------------------------------
\bibliographystyle{IEEEtran}
\bibliography{BLG475E_Phase2_Report}

\end{document}
